\section{Implementation Considerations}

\subsection{Runtime Interface}

The execution runtime exposes a minimal interface:

\begin{quote}
\small
\begin{alltt}
interface ExecutionRuntime \{
    Result execute(CapabilityRequest request)
    List<Capability> listCapabilities()
    HealthStatus healthCheck()
\}
\end{alltt}
\end{quote}

\subsection{Capability Request Format}

Capability requests follow a standardized format:

\begin{quote}
\small
\begin{alltt}
\{
  "capability": "generate\_text",
  "version": "1.0",
  "input": \{
    "prompt": "Explain quantum computing",
    "max\_tokens": 500
  \},
  "context": \{
    "session\_id": "abc123",
    "request\_id": "req456"
  \},
  "provider": \{
    "name": "provider\_a",
    "endpoint": "https://api.provider-a.com"
  \}
\}
\end{alltt}
\end{quote}

\subsection{Result Format}

Results are normalized to a standard structure with execution metrics:

\begin{quote}
\small
\begin{alltt}
\{
  "capability": "generate\_text",
  "status": "success",
  "output": \{
    "text": "Quantum computing is...",
    "tokens\_used": 450
  \},
  "metadata": \{
    "provider": "provider\_a",
    "latency\_ms": 1250,
    "translation\_latency\_ms": 15,
    "retry\_count": 0,
    "timestamp": "2026-03-15T10:30:00Z"
  \}
\}
\end{alltt}
\end{quote}

The \texttt{retry\_count} and \texttt{translation\_latency\_ms} fields provide visibility into execution behavior, enabling the Policy Orchestration Layer to make data-driven routing decisions.
